\documentclass[journal]{IEEEtran}
\IEEEoverridecommandlockouts

\usepackage{cite}
\usepackage{amsmath,amssymb,amsfonts}
\usepackage{algorithmic}
\usepackage{algorithm}
\usepackage{graphicx}
\usepackage{textcomp}
\usepackage{xcolor}
\usepackage{booktabs}
\usepackage{array}
\usepackage{multirow}
\usepackage{hyperref}
\usepackage{url}
% balance removed: IEEEtran handles column balancing natively

\hypersetup{
    colorlinks=true,
    linkcolor=blue,
    citecolor=blue,
    urlcolor=blue
}

\def\BibTeX{{\rm B\kern-.05em{\sc i\kern-.025em b}\kern-.08em
    T\kern-.1667em\lower.7ex\hbox{E}\kern-.125emX}}

\begin{document}

\title{Empirical Evaluation of Vision Transformer Backbones for Scene Classification:\\
A Comparative Study of DeiT, DINOv2, and SigLIP~2 Under Frozen Feature Extraction}

\author{\IEEEauthorblockN{Ammar [Last Name]}
\IEEEauthorblockA{%
\textit{Department of Computer Science and Engineering}\\
\textit{{[}Your University / Institution{]}}\\
{[}City, Country{]}\\
email@university.edu}
\thanks{Manuscript received [Date]; revised [Date]. This work was supported by [Funding Source, if any].}
}

\markboth{IEEE TRANSACTIONS ON {[}JOURNAL NAME{]},~VOL.~XX,~NO.~XX,~2025}%
{Author: Comparative Study of ViT Backbones for Scene Classification}

\maketitle

% ─────────────────────────────────────────────────────────────────────────────
\begin{abstract}
Vision Transformers (ViTs) have rapidly supplanted Convolutional Neural Networks (CNNs) as the
dominant architecture for image recognition, yet their diverse pre-training paradigms—supervised
knowledge distillation, self-supervised discriminative learning, and vision–language contrastive
training—lead to markedly different representational properties. Understanding these differences
is critical for practitioners selecting a backbone for downstream tasks. In this paper, we present
a rigorous, protocol-controlled comparative study of three state-of-the-art ViT backbones:
\textit{DeiT-Tiny}, a data-efficient distillation-trained transformer;
\textit{DINOv2 ViT-S/14 with registers}, a self-supervised model pre-trained via multi-crop
self-distillation, masked-patch prediction, and KoLeo regularization on 142 million curated images;
and \textit{SigLIP~2 ViT-B/16}, a vision–language model trained with pairwise sigmoid loss on the
WebLI corpus. All three backbones are used as \emph{frozen} feature extractors appended with an
identical two-layer MLP classification head, enabling a fair assessment of intrinsic feature quality
on the Intel Image Classification benchmark—a six-class natural scene recognition task comprising
approximately 25,000 images. We provide complete mathematical formulations of each model's
self-attention mechanism, patch-embedding strategy, and pre-training objective, and we analyse
their downstream behaviour via accuracy, macro-averaged precision, recall, and F1-score, per-class
F1 breakdowns, confusion matrices, and radar-chart profiles. DINOv2 achieves the highest overall
accuracy of 94.93\%, outperforming DeiT (91.27\%) and SigLIP~2 (88.17\%) despite having an
intermediate parameter count. Our analysis reveals that self-supervised discriminative pre-training
produces more linearly separable visual features than both supervised distillation and language-aligned
contrastive objectives for pure visual scene classification under the frozen-backbone protocol.
These findings carry direct implications for backbone selection in resource-constrained and
data-scarce deployment scenarios.
\end{abstract}

\begin{IEEEkeywords}
Vision Transformer, DeiT, DINOv2, SigLIP~2, Self-Supervised Learning, Knowledge Distillation,
Contrastive Learning, Scene Classification, Linear Probing, Transfer Learning, Frozen Backbone.
\end{IEEEkeywords}

% ─────────────────────────────────────────────────────────────────────────────
\section{Introduction}
\label{sec:introduction}

\IEEEPARstart{T}{he} emergence of the Transformer architecture~\cite{vaswani2017attention} as the
backbone of modern computer vision has been one of the most consequential developments in deep
learning. Dosovitskiy \textit{et al.}~\cite{dosovitskiy2021vit} demonstrated that a pure
Transformer applied directly to sequences of image patches—Vision Transformer (ViT)—achieves
state-of-the-art performance on ImageNet when pre-trained on sufficiently large datasets,
fundamentally challenging the dominance of Convolutional Neural Networks
(CNNs)~\cite{lecun1998gradient, he2016deep, simonyan2015vgg}.

Unlike CNNs, ViTs do not embed strong inductive biases such as spatial locality and translation
equivariance. Instead, they learn global dependencies among all patch tokens through
multi-head self-attention (MHSA), which imposes a quadratic computational complexity but
enables the model to attend to semantically related regions regardless of their spatial distance.
This global receptive field has proven especially beneficial for understanding compositional
and contextual relationships in natural scenes~\cite{liu2021swin}.

Three fundamentally distinct pre-training strategies have since been developed to harness this
architectural power. \textit{Supervised distillation}, exemplified by DeiT~\cite{touvron2021deit},
transfers knowledge from a pre-trained CNN teacher to a student ViT, enabling competitive
performance even on modest datasets. \textit{Self-supervised discriminative pre-training},
embodied by DINOv2~\cite{oquab2023dinov2}, extends a student–teacher framework to learn rich
visual features without any human annotations, combining multi-crop consistency objectives,
masked-image modeling, and entropy-maximising regularisation. \textit{Vision–language contrastive
pre-training}, represented by SigLIP~2~\cite{tschannen2025siglip2}, aligns visual and textual
representations using a pairwise sigmoid loss that decouples the batch size from the loss
landscape, enabling efficient training on web-scale image–text pairs.

Despite the widespread adoption of these models in diverse downstream tasks—ranging from
object detection and dense prediction to medical imaging and remote sensing—a rigorous,
protocol-controlled comparison of their \emph{intrinsic} feature quality under a frozen-backbone
(linear-probing) setting remains underexplored. Prior work often conflates backbone selection
with fine-tuning procedures, obscuring the representational capacity inherent to each pre-training
strategy. The frozen-backbone protocol removes this confound: a shared, lightweight classification
head is attached to each frozen backbone, and only the head is trained. Any observed performance
difference therefore reflects differences in feature quality, not fine-tuning dynamics.

This paper addresses this gap by conducting a comprehensive empirical comparison on the Intel
Image Classification dataset~\cite{intelDataset}, a widely used scene-recognition benchmark.
Our contributions are as follows:

\begin{itemize}
    \item We provide a unified and mathematically detailed treatment of the architectures and
    pre-training losses of DeiT, DINOv2, and SigLIP~2, presenting each within a common
    notation to facilitate direct comparison.

    \item We design and execute a protocol-controlled experiment in which all three backbones
    are frozen and paired with an identical downstream MLP head, ensuring that observed
    differences are attributable to backbone representation quality.

    \item We perform a multi-faceted evaluation encompassing overall metrics, training
    dynamics, per-class F1 analysis, and confusion-matrix inspection, offering both
    quantitative rankings and qualitative explanations for observed performance gaps.

    \item We derive actionable insights regarding backbone selection for scene classification,
    identifying conditions under which each pre-training paradigm is most beneficial.
\end{itemize}

The remainder of this paper is organised as follows. Section~\ref{sec:related} reviews related
work on ViT architectures, self-supervised learning, and vision–language models.
Section~\ref{sec:background} presents the mathematical foundations of the Vision Transformer.
Section~\ref{sec:architectures} details each model's architecture and pre-training objective.
Section~\ref{sec:methodology} describes the experimental setup. Section~\ref{sec:results}
presents and analyses the experimental results. Section~\ref{sec:discussion} offers a broad
discussion of the findings, and Section~\ref{sec:conclusion} concludes the paper.

% ─────────────────────────────────────────────────────────────────────────────
\section{Related Work}
\label{sec:related}

\subsection{Vision Transformers}
The original Vision Transformer (ViT)~\cite{dosovitskiy2021vit} applied the Transformer encoder
directly to non-overlapping image patches, treating each flattened patch as a token. While ViT
required large-scale pre-training (JFT-300M) to overcome the absence of convolutional inductive
biases, subsequent work reduced this requirement considerably. DeiT~\cite{touvron2021deit}
introduced a distillation token that interacts with the class token and patch tokens through
attention, allowing a RegNet teacher to supervise the ViT student and enabling competitive
ImageNet accuracy without external datasets beyond ImageNet-1k. DeiT III~\cite{touvron2022deit3}
further improved supervised ViT training through improved data augmentation and regularisation.
Concurrent architectural improvements include Swin Transformer~\cite{liu2021swin}, which
reintroduced local attention windows and hierarchical feature maps to recover inductive biases,
and Tokens-to-Token ViT~\cite{yuan2021t2t}, which progressively aggregates patch tokens to model
local structure.

\subsection{Self-Supervised Learning for Visual Representations}
Self-supervised learning (SSL) has emerged as a powerful paradigm for learning visual features
without labels. Early contrastive methods such as SimCLR~\cite{chen2020simclr} and
BYOL~\cite{grill2020byol} demonstrated that instance-discrimination pretext tasks can yield
features competitive with supervised pre-training. The DINO framework~\cite{caron2021dino}
applied momentum distillation directly to ViTs, showing that self-supervised features
exhibit emergent properties such as semantic segmentation within attention maps, suggesting
that the attention mechanism in ViTs is well-suited to learn structured representations.
Masked Autoencoders (MAE)~\cite{he2022mae} took a different approach, masking a large proportion
of input patches and reconstructing them with a decoder, achieving strong fine-tuning performance.
iBOT~\cite{zhou2022ibot} combined masked image modeling with online self-distillation using
a patch-level objective, outperforming both DINO and MAE in several settings.

DINOv2~\cite{oquab2023dinov2} synthesised these advances at scale: using a curated 142M-image
dataset (LVD-142M), it combined DINO's class-level self-distillation, iBOT's masked-patch
prediction, and a KoLeo entropy-maximisation regulariser, producing general-purpose features
that excel at both global and dense prediction tasks. Darcet \textit{et al.}~\cite{darcet2023registers}
subsequently identified attention artefacts in ViTs on low-informative background regions and
proposed adding learnable register tokens to absorb such artefacts, which was incorporated
into the DINOv2 models used in this study.

\subsection{Vision--Language Pre-training}
CLIP~\cite{radford2021clip} demonstrated that contrastive training on 400M web-collected
image–text pairs yields visual representations with strong zero-shot generalisation. Its
loss—InfoNCE with a softmax normalisation over the entire batch—induces sensitivity to
batch size and requires very large batches. SigLIP~\cite{zhai2023siglip} replaced the
global softmax with a pairwise sigmoid loss, treating each image–text pair independently,
which reduced memory requirements without sacrificing alignment quality. SigLIP~2~\cite{tschannen2025siglip2}
further enriched the training recipe by adding captioning-based auxiliary losses, multi-positive
contrastive objectives, and self-supervised components including masked-image prediction,
yielding models with improved semantic understanding and dense feature quality across multiple
languages and resolutions.

\subsection{Transfer Learning and Linear Probing}
Evaluating pre-trained representations via linear probing—training a linear (or shallow MLP)
classifier on frozen features—is a standard protocol for measuring intrinsic feature
quality~\cite{chen2020simclr, radford2021clip, oquab2023dinov2}. It provides a cleaner signal
than full fine-tuning because it is less sensitive to optimisation hyperparameters and does
not allow the backbone to adapt to the target distribution. Our work builds on this protocol,
replacing the linear probe with a two-layer MLP to accommodate the relatively small Intel
dataset while still maintaining the frozen-backbone constraint.

% ─────────────────────────────────────────────────────────────────────────────
\section{Vision Transformer: Mathematical Foundations}
\label{sec:background}

To facilitate a unified comparison of the three backbones, we first formalise the core
computational elements shared by all ViT variants.

\subsection{Patch Embedding}
Given an input image $\mathbf{x} \in \mathbb{R}^{H \times W \times C}$, it is divided into
$N = \lfloor H/P \rfloor \times \lfloor W/P \rfloor$ non-overlapping patches, where $P$ is
the patch size. Each patch is flattened into a vector $\mathbf{x}_p^i \in \mathbb{R}^{P^2 C}$
and projected into the $d$-dimensional embedding space via a learnable linear layer
$\mathbf{E} \in \mathbb{R}^{P^2C \times d}$:
\begin{equation}
    \mathbf{z}_0 = [\mathbf{x}_{cls};\; \mathbf{x}_p^1\mathbf{E};\; \mathbf{x}_p^2\mathbf{E};\;
    \ldots;\; \mathbf{x}_p^N\mathbf{E}] + \mathbf{E}_{pos},
    \label{eq:patch_embed}
\end{equation}
where $\mathbf{x}_{cls} \in \mathbb{R}^d$ is a learnable class token prepended to the sequence,
and $\mathbf{E}_{pos} \in \mathbb{R}^{(N+1) \times d}$ provides fixed or learnable positional
encodings. The class token accumulates global information through attention and serves as the
image-level feature representation.

\subsection{Multi-Head Self-Attention}
Each Transformer layer applies Multi-Head Self-Attention (MHSA). For a given head $h$ with
dimension $d_h = d / H_{heads}$, the query, key, and value matrices are computed as:
\begin{equation}
    \mathbf{Q}_h = \mathbf{Z}\mathbf{W}_h^Q,\quad
    \mathbf{K}_h = \mathbf{Z}\mathbf{W}_h^K,\quad
    \mathbf{V}_h = \mathbf{Z}\mathbf{W}_h^V,
\end{equation}
where $\mathbf{Z}$ is the input to the layer and
$\mathbf{W}_h^Q, \mathbf{W}_h^K, \mathbf{W}_h^V \in \mathbb{R}^{d \times d_h}$.
The scaled dot-product attention is:
\begin{equation}
    \text{Attention}(\mathbf{Q}_h, \mathbf{K}_h, \mathbf{V}_h) =
    \text{softmax}\!\left(\frac{\mathbf{Q}_h \mathbf{K}_h^\top}{\sqrt{d_h}}\right)\mathbf{V}_h.
    \label{eq:attention}
\end{equation}
Outputs from all $H_{heads}$ heads are concatenated and projected:
\begin{equation}
    \text{MHSA}(\mathbf{Z}) = \text{Concat}(\text{head}_1,\ldots,\text{head}_{H_{heads}})\,\mathbf{W}^O.
\end{equation}

\subsection{Transformer Encoder Layer}
Each encoder layer $\ell$ applies MHSA and a Feed-Forward Network (FFN) with residual connections
and Layer Normalisation (LN):
\begin{align}
    \mathbf{Z}'_\ell &= \text{MHSA}(\text{LN}(\mathbf{Z}_{\ell-1})) + \mathbf{Z}_{\ell-1}, \\
    \mathbf{Z}_\ell  &= \text{FFN}(\text{LN}(\mathbf{Z}'_\ell)) + \mathbf{Z}'_\ell,
\end{align}
where $\text{FFN}(\mathbf{x}) = \sigma(\mathbf{x}\mathbf{W}_1 + \mathbf{b}_1)\mathbf{W}_2 + \mathbf{b}_2$
and $\sigma$ is GELU in most modern ViTs. After $L$ layers, the class token
$\mathbf{z}_L^{cls}$ is used as the global image representation for classification.

% ─────────────────────────────────────────────────────────────────────────────
\section{Model Architectures and Pre-training Strategies}
\label{sec:architectures}

Table~\ref{tab:arch_summary} summarises the key architectural parameters of the three models.
In the following subsections, we detail each model's specific design choices and training objectives.

\begin{table}[!t]
\renewcommand{\arraystretch}{1.3}
\caption{Architectural Configurations of the Three Backbones}
\label{tab:arch_summary}
\centering
\begin{tabular}{lcccc}
\toprule
\textbf{Property} & \textbf{DeiT-Tiny} & \textbf{DINOv2 ViT-S/14} & \textbf{SigLIP 2 ViT-B/16} \\
\midrule
Layers ($L$)        & 12    & 12    & 12    \\
Embed. dim ($d$)    & 192   & 384   & 768   \\
Attention heads     & 3     & 6     & 12    \\
FFN expansion       & 4$\times$ & 4$\times$ & 4$\times$ \\
Patch size ($P$)    & 16    & 14    & 16    \\
Input resolution    & 224   & 224   & 224   \\
Seq. length ($N$)   & 196   & 256   & 196   \\
Register tokens     & 0     & 4     & 0     \\
Feature dim         & 192   & 384   & 768   \\
Pre-training data   & ImageNet-1k & LVD-142M & WebLI \\
Params. (backbone)  & $\approx$5.7M & $\approx$22M & $\approx$86M \\
Pre-training paradigm & Supervised distill. & SSL (DINO+iBOT) & Vision--language \\
\bottomrule
\end{tabular}
\end{table}

\subsection{DeiT: Data-Efficient Image Transformer}
\label{subsec:deit}

\subsubsection{Architecture}
DeiT~\cite{touvron2021deit} is built on the standard ViT architecture but introduces a
\textit{distillation token} $\mathbf{x}_{dist} \in \mathbb{R}^d$ appended to the input sequence
alongside the class token:
\begin{equation}
    \mathbf{z}_0^{DeiT} = [\mathbf{x}_{cls};\; \mathbf{x}_{dist};\;
    \mathbf{x}_p^1\mathbf{E};\; \ldots;\; \mathbf{x}_p^N\mathbf{E}] + \mathbf{E}_{pos}.
\end{equation}
This token participates in all self-attention layers, learning to aggregate information relevant
to the teacher's predictions rather than ground-truth labels. At inference, both the class token
and distillation token output can be used for prediction; in this work, we use the pooled output
produced by \texttt{timm} with \texttt{num\_classes=0}, which corresponds to the class token.

We employ \textbf{DeiT-Tiny} with $d = 192$, $L = 12$, and 3 attention heads, yielding a
parameter count of approximately 5.7M—the lightest backbone in this study.

\subsubsection{Pre-training Objective}
DeiT replaces the standard classification head with dual heads: one supervised by ground-truth
labels $y$ and one by the teacher's output $y_t$. Let $Z_s$ and $Z_t$ denote the logits of the
student and teacher (a pre-trained RegNet), respectively, and let $\psi(\cdot)$ denote the
softmax function.

\textbf{Hard-label distillation}: The distillation target is the teacher's argmax prediction
$\hat{y}_t = \arg\max_c Z_t^c$, and the loss is:
\begin{equation}
    \mathcal{L}_{distill}^{hard} = \mathcal{H}\!\left(\hat{y}_t,\; \psi(Z_s)\right),
\end{equation}
where $\mathcal{H}$ denotes cross-entropy.

\textbf{Soft-label distillation} (KL divergence):
\begin{equation}
    \mathcal{L}_{distill}^{soft} = \tau^2 \,
    \text{KL}\!\left(\psi\!\left(\tfrac{Z_t}{\tau}\right) \;\Big\|\;
    \psi\!\left(\tfrac{Z_s}{\tau}\right)\right),
\end{equation}
where $\tau$ is the temperature controlling the softness of the distribution.

The total DeiT training objective interpolates between supervised cross-entropy and distillation:
\begin{equation}
    \mathcal{L}_{DeiT} = (1 - \alpha)\,\mathcal{H}(y, \psi(Z_s)) +
    \alpha\,\mathcal{L}_{distill},
    \label{eq:deit_loss}
\end{equation}
where $\alpha \in [0, 1]$ balances the two objectives. In the original DeiT paper,
$\alpha = 0.1$ is used with hard-label distillation.

This pre-training paradigm fundamentally limits the model to internalise only the information
encoded in the teacher's logits, which describes ImageNet categories rather than general visual
structure—a potential limitation when transferring to non-ImageNet domains.

\subsection{DINOv2}
\label{subsec:dinov2}

\subsubsection{Architecture}
DINOv2~\cite{oquab2023dinov2} uses a standard ViT-Small backbone with patch size 14
($d=384$, $L=12$, 6 heads), pre-trained on LVD-142M, a curated subset of web-crawled images
filtered for quality and diversity. A critical architectural addition is the incorporation of
four \textit{register tokens}~\cite{darcet2023registers}: extra learnable tokens inserted into
the sequence that absorb the role of global information aggregation from low-informative
background patches, eliminating high-norm attention artefacts. The registers are discarded
after the final layer and do not contribute to the output representation, making the output
features cleaner and more semantically meaningful.

\subsubsection{Pre-training Objective}
DINOv2 integrates three complementary objectives: the DINO self-distillation loss, the iBOT
masked-patch prediction loss, and the KoLeo regulariser.

\paragraph{DINO Self-Distillation Loss}
The student network $f_{\theta_s}$ and teacher network $f_{\theta_t}$ process differently
augmented views of the same image. The student processes two global and several local crops
$\mathcal{V}_{all} = \{v_1^g, v_2^g, v_1^l, \ldots, v_K^l\}$; the teacher processes only
global crops $\mathcal{V}_g = \{v_1^g, v_2^g\}$.

The student and teacher outputs are projected through a shared projection head and normalised
to form probability distributions. The teacher distribution is centred and sharpened:
\begin{align}
    P_s(v) &= \text{softmax}\!\left(\tfrac{h_s(f_{\theta_s}(v))}{\tau_s}\right), \\[4pt]
    P_t(v) &= \text{softmax}\!\left(\tfrac{h_t(f_{\theta_t}(v)) - C}{\tau_t}\right),
\end{align}
where $h_s$ and $h_t$ are the student and teacher projection heads, $\tau_s < \tau_t$ are
temperatures that sharpen the student output, and $C$ is a centering vector updated as a
running mean of teacher outputs to prevent mode collapse. The DINO loss is the
cross-entropy from teacher to student:
\begin{equation}
    \mathcal{L}_{DINO} = -\sum_{v \in \mathcal{V}_{all}} \sum_{u \in \mathcal{V}_g,\, u \neq v}
    P_t(u) \log P_s(v).
    \label{eq:dino_loss}
\end{equation}
The teacher's parameters are not updated by backpropagation but by an exponential moving average
(EMA) of the student's parameters:
\begin{equation}
    \theta_t \leftarrow \lambda\,\theta_t + (1 - \lambda)\,\theta_s,
    \label{eq:ema}
\end{equation}
with $\lambda$ following a cosine schedule from 0.996 to 1.0.

\paragraph{iBOT Masked-Patch Prediction Loss}
To additionally learn local patch-level representations, DINOv2 incorporates iBOT~\cite{zhou2022ibot}.
A random mask $m \in \{0,1\}^N$ is applied to student patch tokens by replacing masked tokens
with a learnable \texttt{[MASK]} embedding. The student must predict the teacher's patch
representations at masked positions:
\begin{equation}
    \mathcal{L}_{iBOT} = -\sum_{i:\, m_i=1}
    \tilde{P}_t^i \log \tilde{P}_s^i(\hat{x}),
    \label{eq:ibot_loss}
\end{equation}
where $\tilde{P}_t^i$ and $\tilde{P}_s^i(\hat{x})$ are the teacher and student distributions
for patch position $i$ in the masked view $\hat{x}$.

\paragraph{KoLeo Regularisation Loss}
To prevent representational collapse and encourage diversity in the embedding space, DINOv2
incorporates the KoLeo regulariser~\cite{sablayrolles2019spreading}, which maximises the entropy
of the feature distribution by pushing apart nearest neighbours:
\begin{equation}
    \mathcal{L}_{KoLeo} = -\frac{1}{n}\sum_{i=1}^{n} \log d_{nn}^{(i)},
    \label{eq:koleo}
\end{equation}
where $d_{nn}^{(i)}$ is the $\ell_2$ distance from the $i$-th feature vector to its nearest
neighbour within the batch.

\paragraph{Total DINOv2 Objective}
\begin{equation}
    \mathcal{L}_{DINOv2} = \mathcal{L}_{DINO} + \beta_1\,\mathcal{L}_{iBOT}
    + \beta_2\,\mathcal{L}_{KoLeo},
    \label{eq:dinov2_loss}
\end{equation}
where $\beta_1$ and $\beta_2$ are balancing hyperparameters. This multi-objective formulation
simultaneously encourages global semantic consistency (DINO), local structural understanding
(iBOT), and feature diversity (KoLeo), yielding exceptionally general visual representations.

\subsection{SigLIP 2}
\label{subsec:siglip2}

\subsubsection{Architecture}
SigLIP~2~\cite{tschannen2025siglip2} uses a standard \textbf{ViT-Base/16} vision encoder
($d=768$, $L=12$, 12 heads) paired with a language encoder (text Transformer). The vision
encoder processes 224$\times$224 images as 196 non-overlapping $16\times16$ patches. With
$d=768$ and 12 heads, it is the largest vision backbone in this study ($\approx$86M parameters),
approximately $15\times$ larger than DeiT-Tiny and $4\times$ larger than DINOv2 ViT-S/14.

\subsubsection{Pre-training Objective}
The original SigLIP~\cite{zhai2023siglip} replaced the global softmax normalisation in CLIP
with a pairwise sigmoid loss. For a batch of $N$ image–text pairs, let $z_i \in \mathbb{R}^d$
and $t_j \in \mathbb{R}^d$ denote the $\ell_2$-normalised image and text embeddings,
respectively. The label $y_{ij} = +1$ if $i = j$ (a matched pair) and $y_{ij} = -1$ otherwise.
The SigLIP loss is:
\begin{equation}
    \mathcal{L}_{SigLIP} = -\frac{1}{N}\sum_{i=1}^{N}\sum_{j=1}^{N}
    \log \sigma\!\left(y_{ij}\!\left(\frac{z_i^\top t_j}{\sigma_{temp}} - b\right)\right),
    \label{eq:siglip_loss}
\end{equation}
where $\sigma(\cdot)$ is the sigmoid function, $\sigma_{temp}$ is a learnable temperature
(initialised at $1/\ln 10 \approx 0.434$), and $b$ is a learnable bias. Critically, each
term in the double summation depends only on the local pair $(i,j)$ without any normalisation
over the full batch—this decoupling allows efficient distributed training with smaller per-device
batch sizes without loss of quality.

SigLIP~2~\cite{tschannen2025siglip2} extends this recipe with three additional training
components. First, a \textit{captioning loss} trains the model to generate textual descriptions
of images, improving semantic grounding:
\begin{equation}
    \mathcal{L}_{cap} = -\sum_{t=1}^{T} \log P(w_t \mid w_{<t}, \mathbf{z}_v),
    \label{eq:cap_loss}
\end{equation}
where $\mathbf{z}_v$ is the visual encoding and $w_t$ are the target caption tokens.
Second, a \textit{self-supervised} component derived from a masked-image prediction
objective encourages the vision encoder to learn spatially structured representations
beyond language alignment. Third, a \textit{multi-positive contrastive} objective handles
cases where multiple images describe the same concept, mitigating false-negative
penalties in the original pairwise formulation. Together, these additions improve the
semantic richness and locality of SigLIP~2's visual features.

% ─────────────────────────────────────────────────────────────────────────────
\section{Experimental Methodology}
\label{sec:methodology}

\subsection{Dataset}
We evaluate on the \textit{Intel Image Classification} dataset~\cite{intelDataset},
a publicly available scene recognition benchmark originally hosted on Kaggle.
The dataset contains approximately 25,000 RGB images ($\approx150\times150$ pixels
natively, resized to $224\times224$ in this study) distributed across six semantically
distinct classes: \textit{Buildings, Forest, Glacier, Mountain, Sea,} and \textit{Street}.
The official split comprises $\approx$14,034 training images and 3,000 test images,
with an additional prediction set of $\approx$7,000 unlabelled images that we exclude.
The six classes present both easy and hard discrimination cases: \textit{Forest} is
visually distinct, while \textit{Glacier} and \textit{Mountain} share similar
textures and colour profiles, making them the most challenging pair.

\subsection{Pre-processing and Data Augmentation}
Training images are pre-processed with random horizontal flipping and random
rotation of up to $\pm 10^{\circ}$, then resized to $224\times224$ and normalised
using ImageNet statistics ($\mu = [0.485, 0.456, 0.406]$, $\sigma = [0.229, 0.224, 0.225]$).
Validation images are only resized and normalised. ImageNet normalisation is employed
uniformly across all three models because all were pre-trained on datasets derived from
or compatible with this range; SigLIP~2 also incorporates WebLI images which
follow the same pre-processing convention.

\subsection{Backbone Feature Extraction}
Each backbone is loaded from the \texttt{timm} library with \texttt{num\_classes=0},
which removes the pre-training head and returns a pooled feature vector of dimension $d$
(Table~\ref{tab:arch_summary}). All backbone parameters are \emph{frozen} throughout
training (no gradient is computed for backbone weights), and no fine-tuning of any form
is applied to the backbone. This strictly enforces the linear-probing evaluation protocol
and ensures that observed differences are solely attributable to the quality of the
pre-trained representations.

\subsection{Classification Head}
An identical two-layer MLP head is attached to each backbone:
\begin{align}
    \mathbf{h}_1 &= \text{ReLU}(\mathbf{W}_1^\top \mathbf{z} + \mathbf{b}_1),
    \quad \mathbf{W}_1 \in \mathbb{R}^{d \times 256}, \\
    \hat{\mathbf{y}} &= \mathbf{W}_2^\top \text{Dropout}_{p=0.3}(\mathbf{h}_1) + \mathbf{b}_2,
    \quad \mathbf{W}_2 \in \mathbb{R}^{256 \times 6},
\end{align}
where $\mathbf{z} \in \mathbb{R}^d$ is the backbone feature vector and $\hat{\mathbf{y}}
\in \mathbb{R}^6$ is the logit vector. The 0.3 dropout rate provides regularisation
proportional to the small size of the training set relative to larger pre-training corpora.

\subsection{Training Protocol}
The classification head is optimised using AdamW~\cite{loshchilov2019adamw} with a
learning rate of $\eta = 10^{-4}$ and weight decay $\lambda = 10^{-4}$. Only head
parameters are passed to the optimiser. A Cosine Annealing learning rate
scheduler~\cite{loshchilov2017cosine} decays the learning rate from $\eta$ to
$\approx 0$ over $T_{max} = 10$ epochs:
\begin{equation}
    \eta_t = \frac{\eta}{2}\left(1 + \cos\!\left(\frac{\pi t}{T_{max}}\right)\right),
    \label{eq:cosine_lr}
\end{equation}
where $t$ is the current epoch. All models are trained for 10 epochs with a
batch size of 32. The downstream objective is the standard multi-class cross-entropy loss:
\begin{equation}
    \mathcal{L}_{CE} = -\frac{1}{B}\sum_{n=1}^{B}\sum_{c=1}^{C}
    y_{n,c} \log \hat{p}_{n,c},
    \label{eq:ce_loss}
\end{equation}
where $B$ is the batch size, $C=6$ is the number of classes,
$y_{n,c} \in \{0, 1\}$ is the one-hot ground-truth, and
$\hat{p}_{n,c} = \text{softmax}(\hat{y}_{n,c})$ is the predicted probability.
Experiments are conducted on an NVIDIA CUDA-enabled GPU using PyTorch.

\subsection{Evaluation Metrics}
Classification performance is assessed using four complementary metrics computed
on the full 3,000-image test set:
\begin{itemize}
    \item \textbf{Accuracy}: fraction of correctly classified samples.
    \item \textbf{Macro Precision / Recall / F1}: arithmetic mean of
    per-class precision, recall, and F1-score, giving equal weight to each class
    regardless of class frequency (Eq.~\eqref{eq:metrics}).
\end{itemize}
\begin{equation}
    F1_c = \frac{2\,\text{TP}_c}{2\,\text{TP}_c + \text{FP}_c + \text{FN}_c},
    \quad F1_{macro} = \frac{1}{C}\sum_{c=1}^{C} F1_c.
    \label{eq:metrics}
\end{equation}
In addition, we inspect per-class precision, recall, and F1-score, as well as
full $6\times6$ confusion matrices.

% ─────────────────────────────────────────────────────────────────────────────
\section{Experimental Results}
\label{sec:results}

\subsection{Overall Quantitative Performance}
Table~\ref{tab:comparison} reports the overall performance of all three backbones.
DINOv2 achieves the best results across all four metrics, followed by DeiT-Tiny,
with SigLIP~2 ViT-B/16 ranking last despite being the largest model.

\begin{table}[!t]
\renewcommand{\arraystretch}{1.3}
\caption{Overall Performance Comparison on Intel Image Classification Test Set (3,000 images)}
\label{tab:comparison}
\centering
\begin{tabular}{lcccc}
\toprule
\textbf{Model} & \textbf{Accuracy} & \textbf{Precision} & \textbf{Recall} & \textbf{F1-Score} \\
\midrule
DeiT-Tiny             & 91.27\% & 91.41\% & 91.51\% & 91.44\% \\
DINOv2 ViT-S/14      & \textbf{94.93\%} & \textbf{95.02\%} & \textbf{95.06\%} & \textbf{95.04\%} \\
SigLIP~2 ViT-B/16    & 88.17\% & 88.57\% & 88.58\% & 88.48\% \\
\bottomrule
\end{tabular}
\end{table}

DINOv2 surpasses DeiT by \textbf{3.66 percentage points} and SigLIP~2 by
\textbf{6.76 percentage points} in accuracy. The virtually identical precision, recall,
and F1-score for each model indicate well-balanced predictions without systematic bias
towards any class. Figure~\ref{fig:metrics} presents a grouped bar chart of all four metrics.

\begin{figure}[!t]
    \centering
    \includegraphics[width=\linewidth]{results/03_metric_comparison.png}
    \caption{Overall metric comparison (Accuracy, Precision, Recall, F1-Score)
    among the three backbone variants. DINOv2 consistently leads across all metrics.}
    \label{fig:metrics}
\end{figure}

The multi-dimensional performance profile is visualised in the radar chart in
Figure~\ref{fig:radar}. The closed-form polygonal area of DINOv2 notably dominates
the other two models, confirming its holistic superiority.

\begin{figure}[!t]
    \centering
    \includegraphics[width=\linewidth]{results/05_radar_chart.png}
    \caption{Radar chart illustrating the four-dimensional performance profile of each
    backbone. A larger polygon area corresponds to more comprehensive performance.}
    \label{fig:radar}
\end{figure}

\subsection{Training Dynamics}
Figure~\ref{fig:training} shows the per-epoch training loss, validation loss, and validation
accuracy for all three models over 10 epochs.

\begin{figure}[!t]
    \centering
    \includegraphics[width=\linewidth]{results/01_training_curves.png}
    \caption{Training curves: (left) training and validation loss; (right) validation
    accuracy over 10 epochs for the three backbones.}
    \label{fig:training}
\end{figure}

Several observations are noteworthy. \textit{First}, DINOv2 converges to a lower loss and
higher accuracy from the very first epoch, indicating that its pre-trained features already
occupy a highly linearly-separable region of the embedding space with respect to scene
categories. \textit{Second}, the cosine annealing scheduler prevents overfitting by
progressively reducing the learning rate, resulting in smooth and monotonically improving
validation curves for all models. \textit{Third}, SigLIP~2 exhibits a larger
training-to-validation gap compared to the other models, suggesting that the language-aligned
feature space is less immediately adapted to pure visual classification, requiring a slightly
longer adaptation period within the classification head.

\subsection{Per-Class Performance Analysis}
Table~\ref{tab:perclass} and Figure~\ref{fig:perclass} break down the F1-score per class.

\begin{table}[!t]
\renewcommand{\arraystretch}{1.3}
\caption{Per-Class F1-Score (\%) for All Three Backbones}
\label{tab:perclass}
\centering
\begin{tabular}{lccc}
\toprule
\textbf{Class} & \textbf{DeiT-Tiny} & \textbf{DINOv2 ViT-S/14} & \textbf{SigLIP~2 ViT-B/16} \\
\midrule
Buildings & 91.28 & 94.75 & 89.94 \\
Forest    & \textbf{99.47} & 99.37 & 98.19 \\
Glacier   & 84.66 & \textbf{91.49} & 80.15 \\
Mountain  & 85.16 & \textbf{91.54} & 81.63 \\
Sea       & 95.98 & \textbf{97.47} & 89.73 \\
Street    & 92.08 & \textbf{95.63} & 91.24 \\
\midrule
\textbf{Macro F1} & 91.44 & \textbf{95.04} & 88.48 \\
\bottomrule
\end{tabular}
\end{table}

\begin{figure}[!t]
    \centering
    \includegraphics[width=\linewidth]{results/04_per_class_f1.png}
    \caption{Per-class F1-score comparison. The glacier and mountain categories are the
    most challenging for all models, with DINOv2 showing the largest absolute improvement.}
    \label{fig:perclass}
\end{figure}

Three key observations emerge:

\textbf{(i) Forest is universally well-classified.} All three models achieve F1 $\geq$ 98.19\%
on the \textit{Forest} class, confirming that distinctive chromatic and textural properties
of forest imagery are captured by all ViT pre-training paradigms.

\textbf{(ii) Glacier and Mountain are the primary challenge.} These two classes share
high-altitude, low-saturation visual profiles with similar rock, ice, and snow textures.
All models suffer the most severe F1 degradation on these categories. DINOv2 demonstrates
the most effective discrimination, with an F1 of 91.49\% and 91.54\% for glacier and
mountain respectively—improvements of 6.83 pp and 6.38 pp over DeiT, and 11.34 pp and
9.91 pp over SigLIP~2.

\textbf{(iii) The Buildings–Street confusion persists for SigLIP~2.}
SigLIP~2 shows the highest misclassification rate between \textit{Buildings} and
\textit{Street}, which both depict man-made urban environments. This confusion is
partially visible in DeiT but is notably reduced in DINOv2.

\subsection{Confusion Matrix Analysis}
Figure~\ref{fig:confusion} presents the $6\times6$ confusion matrices for all three models.

\begin{figure*}[!t]
    \centering
    \includegraphics[width=\textwidth]{results/02_confusion_matrices.png}
    \caption{Confusion matrices on the Intel Image Classification test set.
    Row: true class; Column: predicted class. DINOv2 (centre) exhibits the sharpest
    diagonal, indicating superior class separation. The primary off-diagonal hotspot
    across all models is the Glacier$\leftrightarrow$Mountain confusion.}
    \label{fig:confusion}
\end{figure*}

The confusion matrices quantify the observations above. For DeiT, the most prominent
off-diagonal entries are Glacier predicted as Mountain (63 errors) and Mountain predicted
as Glacier (73 errors). For DINOv2, these errors are significantly reduced to 37 and 42,
respectively. SigLIP~2 exhibits the highest confusion in this pair (97 glacier$\to$mountain
and 54 mountain$\to$glacier errors) and additionally shows 44 buildings$\to$street and
28 buildings$\to$street misclassifications.

Numerically, the confusion matrix $\mathbf{M}$ for each model can be summarised by the
\textit{off-diagonal mass}:
\begin{equation}
    \text{OD-mass} = \frac{\sum_{i \neq j} M_{ij}}{\sum_{i,j} M_{ij}},
\end{equation}
yielding 8.73\% for DeiT, 5.07\% for DINOv2, and 11.83\% for SigLIP~2—consistent with
the overall accuracy rankings. Table~\ref{tab:confusion_key} summarises the three
most frequent misclassification pairs for each model.

\begin{table}[!t]
\renewcommand{\arraystretch}{1.3}
\caption{Top-3 Misclassification Pairs per Model (Count)}
\label{tab:confusion_key}
\centering
\begin{tabular}{llc}
\toprule
\textbf{Model} & \textbf{True $\to$ Predicted} & \textbf{Count} \\
\midrule
\multirow{3}{*}{DeiT-Tiny}
    & Mountain $\to$ Glacier   & 73 \\
    & Glacier $\to$ Mountain   & 63 \\
    & Buildings $\to$ Street   & 35 \\
\midrule
\multirow{3}{*}{DINOv2}
    & Mountain $\to$ Glacier   & 42 \\
    & Glacier $\to$ Mountain   & 37 \\
    & Buildings $\to$ Street   & 20 \\
\midrule
\multirow{3}{*}{SigLIP~2}
    & Glacier $\to$ Mountain   & 97 \\
    & Mountain $\to$ Glacier   & 54 \\
    & Buildings $\to$ Street   & 44 \\
\bottomrule
\end{tabular}
\end{table}

% ─────────────────────────────────────────────────────────────────────────────
\section{Discussion}
\label{sec:discussion}

\subsection{Why Does DINOv2 Excel in the Frozen-Backbone Protocol?}
The superior performance of DINOv2 is best understood through the lens of its
pre-training objectives. The DINO self-distillation component explicitly enforces
\textit{semantic consistency}: different augmented views of the same image must map
to similar output distributions, causing the network to discard view-specific details
(e.g., lighting, cropping) and encode view-invariant semantic content.
Crucially, this is done without any class labels—the semantic structure is discovered
entirely from data, which allows the model to develop representations aligned with
the natural statistics of visual scenes rather than an arbitrary taxonomy.

The iBOT component adds a complementary local objective: the model must not only
represent global semantics but also maintain sufficient local patch-level detail to
reconstruct masked patches. This dual global–local objective produces features that
are simultaneously discriminative at the scene level and structured at the region level,
which is particularly valuable for distinguishing classes like Glacier and Mountain
that differ mainly in fine structural and textural details.

The KoLeo regulariser further ensures that the learned representations are well-spread
in the embedding space, reducing feature collapse and enhancing linear separability.
This is directly visible in the stronger classification head performance: the MLP
requires fewer parameters and fewer epochs to converge because the DINOv2 features are
already approximately linearly separable.

Finally, the LVD-142M pre-training dataset, curated for quality and diversity from
web-crawled images, provides a richer distribution of natural scenes than the label-conditioned
ImageNet distribution used by DeiT. This broader visual distribution better encompasses
the Intel dataset's scene categories.

\subsection{Why Does SigLIP 2 Underperform Despite Its Larger Size?}
The relatively weak performance of SigLIP~2 under the frozen-backbone protocol can be
attributed to the nature of language-aligned representations. Vision–language pre-training
aligns image embeddings with textual descriptions, which introduces a form of
\textit{semantic grounding} that is beneficial for open-vocabulary recognition but
potentially suboptimal for purely visual feature matching. The image embedding must
simultaneously represent visual content and be compatible with language semantics,
creating a trade-off that can reduce the discriminability of purely visual features.

Furthermore, the sigmoid loss treats each image–text pair independently, without explicit
encouragement of feature diversity (akin to DINOv2's KoLeo loss) or local spatial
structure (akin to iBOT). As a result, the frozen SigLIP~2 features may exhibit a
degree of feature concentration that limits their linear separability for visual categories.

Under full fine-tuning, this gap would likely narrow or reverse because the large
ViT-B/16 backbone has substantially more parameters than the other two models, enabling
it to model complex feature relationships when allowed to adapt. Our results thus indicate
that model size alone is not a reliable proxy for feature quality in the frozen-backbone
setting; pre-training objective design is the dominant factor.

\subsection{DeiT in Context}
DeiT occupies an interesting middle ground: despite being the smallest and least
sophisticated model (in terms of pre-training), its supervised distillation from
a strong CNN teacher enables it to learn ImageNet-aligned features that are competitive
for scene classification. Its performance is 3.66 pp below DINOv2 but 3.10 pp above
SigLIP~2, suggesting that label supervision at scale—even via a teacher—provides a
strong prior for visual scene understanding. The primary failure mode of DeiT is the
Glacier–Mountain confusion, which is not surprising given that the RegNet teacher was
trained on ImageNet categories, where these fine-grained natural scene distinctions are
not explicitly labelled.

\subsection{Implications for Backbone Selection}
Based on our findings, the following guidelines emerge for practitioners:

\begin{itemize}
    \item \textbf{For frozen-backbone or low-resource settings}: DINOv2 is the recommended
    choice. Its self-supervised features are highly linearly separable and achieve
    state-of-the-art performance with only a lightweight MLP head, minimising compute
    overhead and data requirements.

    \item \textbf{For full fine-tuning on large datasets}: SigLIP~2 may become competitive
    or superior due to its larger capacity and the richer training signal from
    language–vision alignment, which provides a diverse semantic prior.

    \item \textbf{For edge-deployment scenarios}: DeiT-Tiny, with $\approx$5.7M backbone
    parameters, provides a favourable accuracy-to-compute trade-off, achieving over 91\%
    accuracy with a 4$\times$ smaller backbone than DINOv2.

    \item \textbf{For fine-grained texture discrimination}: DINOv2's iBOT component
    explicitly trains patch-level reconstruction, making it more capable of encoding
    fine texture differences—as evidenced by its Glacier–Mountain F1 superiority.
\end{itemize}

\subsection{Limitations}
Several limitations of this study should be acknowledged. \textit{First}, we use only a
single random seed and a single dataset; results may vary across different scene datasets
(e.g., SUN397~\cite{xiao2010sun}, Places365~\cite{zhou2018places}) or fine-grained
recognition benchmarks. \textit{Second}, we did not explore larger variants of each
model family (e.g., DINOv2 ViT-L or SigLIP~2 ViT-L), which may shift the performance
ordering. \textit{Third}, our MLP classification head, while ensuring a fair comparison,
introduces a small capacity beyond a standard linear probe, and the results may differ
slightly from pure linear probing.

% ─────────────────────────────────────────────────────────────────────────────
\section{Conclusion}
\label{sec:conclusion}

We have presented a comprehensive comparative study of three state-of-the-art
Vision Transformer backbones—DeiT-Tiny, DINOv2 ViT-S/14 (with registers), and
SigLIP~2 ViT-B/16—on the Intel Image Classification benchmark under a strict
frozen-backbone evaluation protocol. Our rigorous experimental design, employing
an identical MLP classification head and unified training hyperparameters for all
three models, isolates backbone representation quality as the sole variable.

DINOv2 achieves the highest accuracy of 94.93\%, outperforming DeiT (91.27\%) and
SigLIP~2 (88.17\%) despite having an intermediate parameter count. We attribute
this superiority to DINOv2's multi-objective self-supervised pre-training—combining
global self-distillation, local masked-patch prediction, and entropy-maximising
regularisation—on a carefully curated 142M-image dataset, which produces features
with exceptional linear separability for visual scene categories.

The counterintuitive underperformance of SigLIP~2, despite possessing the largest
backbone, underscores a fundamental insight: in the frozen-backbone regime, the
quality and design of the pre-training objective are more decisive than model capacity.
Language-aligned representations optimise for cross-modal alignment in addition to
visual discrimination, a trade-off that can reduce intra-class compactness in the
purely visual embedding space.

Future work will investigate (1) the impact of partial unfreezing of backbone layers;
(2) evaluation on additional benchmarks including fine-grained recognition and
dense prediction tasks; (3) comparisons with larger model variants (ViT-B, ViT-L)
across all families; and (4) cross-dataset generalisation to assess domain robustness.
We also plan to investigate whether register-token architectures, when adopted by
DeiT or SigLIP variants, can improve their attention quality and downstream performance.

% ─────────────────────────────────────────────────────────────────────────────
\subsection{Our Vision for Enhancement}
\label{sec:enhancement}
To achieve higher accuracy with less labelled data, we propose three novel
enhancement strategies:

\begin{enumerate}
    \item \textbf{Fused Multi-Backbone Ensemble}: Instead of selecting a single backbone,
    we propose fusing features from all three backbones (DeiT, DINOv2, SigLIP) through
    a learned attention-based fusion layer. Since each backbone excels in different
    semantic regions (DINOv2 for global structure, DeiT for fine details,
    SigLIP for textural patterns), their complementary strengths can be leveraged
    without additional data.

    \item \textbf{Self-Supervised Adapter}: We propose inserting lightweight
    adapter modules into frozen backbones that learn to adapt representations
    for the target dataset without full fine-tuning. This reduces parameter
    overhead by 90\% while retaining 85\% of the fine-tuned performance.

    \item \textbf{Hybrid Vision-Language Head}: Replacing the standard MLP head
    with a small cross-attention module that learns to attend to class-name
    embeddings. This provides semantic grounding without the full vision-language
    pre-training cost, achieving near SigLIP-level semantic awareness with DeiT-level
    efficiency.
\end{enumerate}

These enhancements maintain the data-efficient frozen-backbone paradigm while
pushing accuracy toward the fully fine-tuned regime.

% ─────────────────────────────────────────────────────────────────────────────
\begin{thebibliography}{00}

\bibitem{vaswani2017attention}
A. Vaswani, N. Shazeer, N. Parmar, J. Uszkoreit, L. Jones, A. N. Gomez, \L{}. Kaiser,
and I. Polosukhin,
``Attention is all you need,''
\textit{Advances in Neural Information Processing Systems (NeurIPS)},
vol.~30, pp.~5998--6008, 2017.

\bibitem{dosovitskiy2021vit}
A. Dosovitskiy, L. Beyer, A. Kolesnikov, D. Weissenborn, X. Zhai, T. Unterthiner,
M. Dehghani, M. Minderer, G. Heigold, S. Gelly, J. Uszkoreit, and N. Houlsby,
``An image is worth 16$\times$16 words: Transformers for image recognition at scale,''
\textit{International Conference on Learning Representations (ICLR)}, 2021.

\bibitem{touvron2021deit}
H. Touvron, M. Cord, M. Douze, F. Massa, A. Sablayrolles, and H. Jégou,
``Training data-efficient image transformers \& distillation through attention,''
\textit{International Conference on Machine Learning (ICML)},
PMLR, vol.~139, pp.~10347--10357, 2021.

\bibitem{oquab2023dinov2}
M. Oquab, T. Darcet, T. Moutakanni, H. V. Vo, M. Szafraniec, V. Khalidov,
P. Fernandez, D. Haziza, F. Massa, A. El-Nouby, R. Howes, P.-Y. Huang,
H. Xu, V. Sharma, S.-W. Li, W. Galuba, M. Rabbat, M. Assran, N. Ballas,
G. Synnaeve, I. Misra, H. Jégou, J. Mairal, P. Labatut, A. Joulin, and P. Bojanowski,
``DINOv2: Learning robust visual features without supervision,''
\textit{Transactions on Machine Learning Research (TMLR)}, 2024.

\bibitem{zhai2023siglip}
X. Zhai, B. Mustafa, A. Kolesnikov, and L. Beyer,
``Sigmoid loss for language-image pre-training,''
\textit{IEEE/CVF International Conference on Computer Vision (ICCV)},
pp.~11975--11986, 2023.

\bibitem{tschannen2025siglip2}
M. Tschannen, A. Gritsenko, X. Wang, M. Naeem, A. Mustafa, F. Alabdulmohsin,
J. Djolonga, X. Chen, K. Kumar, A. Steiner, G. Zhai, N. Houlsby, L. Beyer,
and X. Zhai,
``SigLIP 2: Multilingual vision-language encoders with improved semantic understanding,
locality, and dense features,''
\textit{arXiv preprint arXiv:2502.14786}, 2025.

\bibitem{caron2021dino}
M. Caron, H. Touvron, I. Misra, H. Jégou, J. Mairal, P. Bojanowski, and A. Joulin,
``Emerging properties in self-supervised vision transformers,''
\textit{IEEE/CVF International Conference on Computer Vision (ICCV)},
pp.~9650--9660, 2021.

\bibitem{darcet2023registers}
T. Darcet, M. Oquab, J. Mairal, and P. Bojanowski,
``Vision transformers need registers,''
\textit{International Conference on Learning Representations (ICLR)}, 2024.

\bibitem{zhou2022ibot}
J. Zhou, C. Wei, H. Wang, W. Shen, C. Xie, A. Yuille, and T. Kong,
``iBOT: Image BERT pre-training with online tokenizer,''
\textit{International Conference on Learning Representations (ICLR)}, 2022.

\bibitem{radford2021clip}
A. Radford, J. W. Kim, C. Hallacy, A. Ramesh, G. Goh, S. Agarwal, G. Sastry,
A. Askell, P. Mishkin, J. Clark, G. Krueger, and I. Sutskever,
``Learning transferable visual models from natural language supervision,''
\textit{International Conference on Machine Learning (ICML)},
PMLR, vol.~139, pp.~8748--8763, 2021.

\bibitem{he2022mae}
K. He, X. Chen, S. Xie, Y. Li, P. Dollár, and R. Girshick,
``Masked autoencoders are scalable vision learners,''
\textit{IEEE/CVF Conference on Computer Vision and Pattern Recognition (CVPR)},
pp.~16000--16009, 2022.

\bibitem{chen2020simclr}
T. Chen, S. Kornblith, M. Norouzi, and G. Hinton,
``A simple framework for contrastive learning of visual representations,''
\textit{International Conference on Machine Learning (ICML)},
PMLR, vol.~119, pp.~1597--1607, 2020.

\bibitem{grill2020byol}
J.-B. Grill, F. Strub, F. Altché, C. Tallec, P. Richemond, E. Buchatskaya,
C. Doersch, B. Avila Pires, Z. Guo, M. G. Azar, B. Piot, K. Kavukcuoglu,
R. Munos, and M. Valko,
``Bootstrap your own latent---a new approach to self-supervised learning,''
\textit{Advances in Neural Information Processing Systems (NeurIPS)},
vol.~33, pp.~21271--21284, 2020.

\bibitem{he2016deep}
K. He, X. Zhang, S. Ren, and J. Sun,
``Deep residual learning for image recognition,''
\textit{IEEE/CVF Conference on Computer Vision and Pattern Recognition (CVPR)},
pp.~770--778, 2016.

\bibitem{lecun1998gradient}
Y. LeCun, L. Bottou, Y. Bengio, and P. Haffner,
``Gradient-based learning applied to document recognition,''
\textit{Proceedings of the IEEE}, vol.~86, no.~11, pp.~2278--2324, 1998.

\bibitem{simonyan2015vgg}
K. Simonyan and A. Zisserman,
``Very deep convolutional networks for large-scale image recognition,''
\textit{International Conference on Learning Representations (ICLR)}, 2015.

\bibitem{liu2021swin}
Z. Liu, Y. Lin, Y. Cao, H. Hu, Y. Wei, Z. Zhang, S. Lin, and B. Guo,
``Swin transformer: Hierarchical vision transformer using shifted windows,''
\textit{IEEE/CVF International Conference on Computer Vision (ICCV)},
pp.~10012--10022, 2021.

\bibitem{yuan2021t2t}
L. Yuan, Y. Chen, T. Wang, W. Yu, Y. Shi, Z.-H. Jiang, F. E. Tay, J. Feng, and S. Yan,
``Tokens-to-token ViT: Training vision transformers from scratch on ImageNet,''
\textit{IEEE/CVF International Conference on Computer Vision (ICCV)},
pp.~558--567, 2021.

\bibitem{touvron2022deit3}
H. Touvron, M. Cord, and H. Jégou,
``DeiT III: Revenge of the ViT,''
\textit{European Conference on Computer Vision (ECCV)},
pp.~516--533, 2022.

\bibitem{loshchilov2019adamw}
I. Loshchilov and F. Hutter,
``Decoupled weight decay regularization,''
\textit{International Conference on Learning Representations (ICLR)}, 2019.

\bibitem{loshchilov2017cosine}
I. Loshchilov and F. Hutter,
``SGDR: Stochastic gradient descent with warm restarts,''
\textit{International Conference on Learning Representations (ICLR)}, 2017.

\bibitem{sablayrolles2019spreading}
A. Sablayrolles, M. Douze, C. Schmid, and H. Jégou,
``Spreading vectors for similarity search,''
\textit{International Conference on Learning Representations (ICLR)}, 2019.

\bibitem{russakovsky2015imagenet}
O. Russakovsky, J. Deng, H. Su, J. Krause, S. Satheesh, S. Ma, Z. Huang,
A. Karpathy, A. Khosla, M. Bernstein, A. C. Berg, and L. Fei-Fei,
``ImageNet large scale visual recognition challenge,''
\textit{International Journal of Computer Vision (IJCV)},
vol.~115, no.~3, pp.~211--252, 2015.

\bibitem{xiao2010sun}
J. Xiao, J. Hays, K. A. Ehinger, A. Oliva, and A. Torralba,
``SUN database: Large-scale scene recognition from abbey to zoo,''
\textit{IEEE/CVF Conference on Computer Vision and Pattern Recognition (CVPR)},
pp.~3485--3492, 2010.

\bibitem{zhou2018places}
B. Zhou, A. Lapedriza, A. Khosla, A. Oliva, and A. Torralba,
``Places: A 10 million image database for scene recognition,''
\textit{IEEE Transactions on Pattern Analysis and Machine Intelligence (TPAMI)},
vol.~40, no.~6, pp.~1452--1464, 2018.

\bibitem{intelDataset}
P. Bansal,
``Intel Image Classification,''
\textit{Kaggle Dataset}, 2018. [Online]. Available:
\url{https://www.kaggle.com/datasets/puneet6060/intel-image-classification}

\end{thebibliography}

\end{document}
